\paragraph{Leakage.}
With a random, sample-level split of the same size, in which every signer appears in both training
and test data, the Conv-Transformer reaches \numLeakRandom{}\,\% test top-1 instead of
\numLeakSigner{}\,\%. A random split thus overestimates the accuracy on new users by \numLeakGap{}
points, more than the effect of any architectural or preprocessing choice studied here.

\paragraph{Ablations.}
\Cref{fig:ablations} changes one factor at a time with respect to the full model
(\numRefValTop{}\,\% validation top-1; seed-to-seed standard deviation of 0.1 point). Early stopping
proved to be a confounding factor. Two variants initially stopped on a validation plateau after 21
and 31 epochs, while the learning rate was still high, and appeared to lose 4.4 and 2.7 points.
Re-run over the full 50-epoch schedule, their effects reduce to \numAblNoHandlocal{} (no hand-local
features) and \numAblXyz{} (depth added) points. The other variants marked with $\dagger$
(\numAblEarlyStopped{} in total) stopped between epochs 41 and 46; their effects may be slightly
overestimated, and our compute budget did not allow re-running them.

Under this caveat, the canonical dominant hand has the largest effect (\numAblNoCanonical{} points
without it): otherwise, the same sign occupies either hand slot depending on the handedness of the
signer and on mirroring. Velocities (\numAblNoVelocity{}), augmentation (\numAblNoAugment{}), depth
(\numAblXyz{}, consistent with the noise of monocular depth estimates) and the hand-local shape
features (\numAblNoHandlocal{}) have moderate effects.

Three findings depart from common practice. First, the lips and the upper body bring no
measurable gain on this vocabulary (\numAblNoLips{} without lips, \numAblHandsOnly{} with hands
only): the 250 signs are mainly distinguished by the hands, and mouthing is presumably rare in
dictionary-like recordings. Second, shorter sequences perform as well ($T'=32$: \numAblTThreeTwo{})
and longer ones no better ($T'=96$: \numAblTNineSix{}), the median clip containing only 22 frames.
Third, zero-padding with a mask outperforms resampling (\numAblPad{}): resampling stretches short
clips and removes the duration cue, whereas padding preserves the original temporal dynamics. The
resampled configuration, fixed before the ablations, is kept in the rest of the study, but padding
is the better choice for future work.

\paragraph{Training budget and ensembles.}
Training the Conv-Transformer for 150 instead of 50 epochs does not help (\numLongVal{}\,\% against
\numRefValTop{}\,\% on validation): the late rise of the curves in \cref{fig:curves} reflects the
annealing of the learning rate, not under-training. Averaging the predictions of the three seeds,
by contrast, improves test top-1 to \numAslEnsTop{}\,\% on ASL (\numTrTestTopMean{}\,\% for a
single model) and to \numLsfbEnsTop{}\,\% on LSFB (\numLsfbScratchTest{}\,\%), at three times the
inference cost.

\begin{figure}[t]
  \centering
  \includegraphics[width=0.62\linewidth]{ablations}
  \caption{One-factor ablations of the Conv-Transformer on the ASL validation signers (seed 0).
  $\dagger$: early stopping ended the run before the end of the 50-epoch schedule (epochs 41--46),
  so the effect may be overestimated.}
  \label{fig:ablations}
\end{figure}
